\documentclass[11pt,a4paper]{article}
\usepackage[utf8]{inputenc}
\usepackage[T1]{fontenc}
\usepackage{lmodern}
\usepackage{amsmath,amssymb}
\usepackage[margin=28mm]{geometry}
\usepackage{microtype}
\usepackage{hyperref}
\hypersetup{colorlinks=true,linkcolor=blue,urlcolor=blue}
\setlength{\emergencystretch}{2em}
\setcounter{secnumdepth}{0}
\providecommand{\tightlist}{\setlength{\itemsep}{0pt}\setlength{\parskip}{0pt}}
\title{Removing the integral-Newton restriction for GL\_n and Sp\_2n}
\author{}
\date{}
\begin{document}
\maketitle
This argument concerns GL\_n and Sp\_2n in the equal- and
mixed-characteristic settings of the uploaded paper. It does not claim
the extension for every split reductive group, nor a Lean formalization.
No claim of novelty is made.

\textbf{Further extension:} \url{general-extension.md}, using the
additional \href{gsp-proof.md}{GSp argument}, treats type-A/C Newton
Levis and split G2. The conjecture for arbitrary split reductive groups
is still not proved.

\section{The statement}\label{the-statement}

Let k be algebraically closed and take either L = k((t)), O =
k{[}{[}t{]}{]}, or L = W(k){[}1/p{]}, O = W(k), t = p, as in §1.1 of the
uploaded paper. Assume the characteristic hypotheses of Xuhua He,
\emph{Irreducible components of affine Lusztig varieties},
arXiv:2502.16441v1: char(k) = 0 or char(k) does not divide the order of
the Weyl group. Let

\[
 G=\mathrm{GL}_n\quad\text{or}\quad G=\mathrm{Sp}_{2n}.
\]

Let gamma in G(L) be regular semisimple. Let mu be dominant and
Y\_mu(gamma) nonempty. Associate {[}b{]} to gamma as in He, §1.2, and
let lambda\_gamma = lambda\_b be the best integral approximation defined
in He, §3.2. Then

\[
 \boxed{\#\bigl(Z_{G(L)}(\gamma)\backslash
 \Sigma^{\mathrm{top}}(Y_\mu(\gamma))\bigr)
 =\dim V_\mu(\lambda_\gamma).}
\]

There is no integrality assumption on the Newton point. All numbered
references to propositions and theorems below refer to the uploaded
paper unless stated otherwise.

\section{1. The component-count
reduction}\label{the-component-count-reduction}

Put w = w\_0 t\^{}mu and fix a reduction tree T. Let R be the paths
corresponding to gamma that contribute to the maximal dimension,
equivalently those satisfying

\[
 \ell(p)+\ell(\operatorname{end}(p))
 -\langle\nu_\gamma,2\rho\rangle=\dim X_w(b).
\]

Theorem 1.4 identifies these paths with the top-contributing paths on
the affine Lusztig side. The component bijections used in §3.6, Theorem
2.8, Proposition 2.5, and Theorem 3.2(1) give, without requiring
integrality,

\[
 C_\mu(\gamma)
 :=\#\bigl(Z_{G(L)}(\gamma)\backslash
 \Sigma^{\mathrm{top}}(Y_\mu(\gamma))\bigr)
 =\sum_{p\in\mathcal R}n_{p,\gamma},
 \qquad
 |\mathcal R|=\dim V_\mu(\lambda_\gamma).
\]

Thus it suffices to prove n\_\{p,gamma\} = 1 for every p in R.

Every such path has a very special endpoint. One can use the stabilizer
argument in §3.2 directly, so no separate virtual-dimension assertion is
needed here. Indeed, X\_p(b) is J\_b-stable and has the same dimension
as X\_w(b), giving an equivariant injection of its top components into
those of X\_w(b). Its fibration over X\_end(p)(b) gives an equivariant
bijection on top components. The stabilizers are therefore the
stabilizers of top components of X\_w(b), which are very special by
Theorem 3.2(2) and the component bijection between X\_w(b) and X\_mu(b).
As in §3.2, this means that the parahoric P\_p has maximal log-volume
among the tau-stable parahorics in the Newton Levi M\_p.

Write end(p) = u tau as in §2.5. Choose the corresponding conjugate
gamma\textquotesingle{} in P\_p dot(tau). Proposition 2.7 and its proof
identify the top components of the partial affine Springer fiber with
the top components of its regular locus. Consequently it is enough to
show

\[
 Z_{M_p(L)}(\gamma')
 \text{ acts transitively on }
 Y^{M_p,\mathrm{reg}}_{K_p,\tau}(\gamma').
 \tag{1}
\]

Theorem 1.3 identifies this centralizer with the original centralizer of
gamma.

\section{2. Very special parahorics for a general-linear slope
block}\label{very-special-parahorics-for-a-general-linear-slope-block}

Consider GL\_N, where N = md, with an isoclinic Newton slope r/d in
lowest terms. The determinant valuation is mr, and the length-zero
element is tau = omega\^{}\{mr\}, where omega is the usual length-zero
rotation with determinant valuation one.

For N \textgreater= 2, label the affine A\_\{N-1\} vertices cyclically
by Z/NZ. Conjugation by tau rotates them by mr. Since gcd(r,d) = 1, its
orbits are exactly the d-element congruence classes modulo m.

A tau-stable spherical subset K is the complement of a nonempty union of
these orbits. If q congruence classes are deleted, let a\_1,...,a\_q be
the cyclic gaps between their representatives in one period of length m.
Then a\_i \textgreater= 1, their sum is m, and the log-volume is

\[
 n_K=d\sum_{i=1}^q\frac{a_i(a_i-1)}2.
\]

For q \textgreater{} 1, this is strictly less than d m(m-1)/2. Hence
maximal log-volume occurs exactly when a single orbit is deleted. For m
= 1 this gives K empty. For d = 1 it gives an ordinary hyperspecial
parahoric. GL\_1 is the root-free version of the same lattice
description.

Accordingly, a very special P\_K in this block stabilizes a partial
periodic lattice chain

\[
 \cdots\supset\Lambda_i\supset\Lambda_{i+1}\supset\cdots,
 \qquad\Lambda_{i+d}=t\Lambda_i,
 \qquad\dim_k(\Lambda_i/\Lambda_{i+1})=m.
 \tag{2}
\]

The coset condition defining Y\_\{K,tau\}(gamma) is

\[
 \gamma\Lambda_i=\Lambda_{i+r}\quad(i\in\mathbb Z).
 \tag{3}
\]

This works for every omitted orbit: choosing a different orbit changes
the reference partial chain but not (2)-\/-(3).

\section{3. The integral root
construction}\label{the-integral-root-construction}

We prove the transitivity in (1) for an arbitrary isoclinic GL\_N block,
with no coprimality condition on the determinant valuation and N.

Let gamma be regular semisimple on V = L\^{}N, all of whose eigenvalues
have valuation r/d. Set

\[
 \beta=t^{-r}\gamma^d,
 \qquad B=O[\beta].
\]

All eigenvalues of beta are units. Its minimal polynomial has
coefficients in O and unit constant coefficient. Thus B is finite over
O, beta is a unit of B, and B is t-adically complete. It is a
commutative subalgebra of L{[}gamma{]}.

There is an element h in B\^{}times such that

\[
 h^d=\beta.
 \tag{4}
\]

To see this, B/tB is a finite Artinian k-algebra. Its residue fields are
k and beta is a unit in each one. Since k is algebraically closed and d
is invertible in k, choose a d-th root in each residue field and lift it
through the nilpotent radical. The derivative d h\^{}\{d-1\} is a unit.
Hensel lifting in the complete ring B then gives (4). Here d is
invertible because d \textless= N and the paper\textquotesingle s GL\_N
characteristic hypothesis excludes all primes \textless= N.

Choose integers a,b with ad + br = 1 and define the fixed operator

\[
 \pi=t^a\gamma^b h^{-b}.
 \tag{5}
\]

All factors commute. Equations (4)-\/-(5) give the exact identities

\[
 \boxed{\pi^d=t,\qquad\gamma=\pi^r h.}
 \tag{6}
\]

For example, pi\^{}d = t\^{}\{ad\} gamma\^{}\{bd\} beta\^{}\{-b\} =
t\^{}\{ad+br\}, and pi\^{}r h = t\^{}\{ar\} gamma\^{}\{br\} h\^{}\{ad\}
= gamma\^{}\{br+ad\}.

Since X\^{}d-t is Eisenstein, F = L{[}pi{]} is a totally ramified field
of degree d and O\_F = O{[}pi{]}. In equal characteristic F is k((pi)).
The vector space V has dimension m over F, and h is F-linear. The
argument below is linear algebra over O\_F, so it does not require
applying a geometric theorem over this auxiliary, possibly ramified,
field.

Moreover, putting A = L{[}gamma{]}, one has

\[
 A=F[h].
 \tag{7}
\]

Both inclusions follow from (5)-\/-(6) and h in B. Because gamma is
regular semisimple, A is an etale L-algebra of dimension N, acting with
a cyclic vector on V. Equation (7) implies that h is regular semisimple
on the m-dimensional F-vector space V. The centralizers agree exactly:

\[
 Z_{\mathrm{GL}_N(L)}(\gamma)
 =A^\times
 =Z_{\mathrm{GL}_m(F)}(h).
 \tag{8}
\]

Equivalently, an L-linear operator commuting with gamma commutes with
beta, h, and pi; conversely an F-linear operator commuting with h
commutes with gamma.

\section{4. Identifying the lattice problem and its regular
locus}\label{identifying-the-lattice-problem-and-its-regular-locus}

Suppose a partial chain satisfies (2)-\/-(3). Then

\[
 \beta\Lambda_i=t^{-r}\gamma^d\Lambda_i=\Lambda_i.
\]

It follows that B preserves Lambda\_i. Both h and h\^{}\{-1\} belong to
B, so h Lambda\_i = Lambda\_i. Using (5) and ad+br = 1,

\[
 \pi\Lambda_i=\Lambda_{i+ad+br}=\Lambda_{i+1}.
 \tag{9}
\]

Consequently Lambda\_0 is an O\_F-lattice of rank m, and

\[
 \Lambda_i=\pi^i\Lambda_0,
 \qquad h\Lambda_0=\Lambda_0.
 \tag{10}
\]

Conversely, any h-stable O\_F-lattice Lambda gives the partial chain
pi\^{}i Lambda. This chain has the period and quotient dimensions in
(2), and gamma = pi\^{}r h gives (3). This is an equivariant bijection
of lattice points, with the centralizers identified by (8). Importantly,
F and h were chosen from the fixed gamma, independently of the lattice.

The regular loci also match. The reductive quotient of P\_K is
(GL\_m)\^{}d. The automorphism induced by tau permutes these d factors
cyclically, since gcd(r,d) = 1. For an element in this cyclic coset, its
centralizer is isomorphic to the centralizer of the product around the
cycle in GL\_m: the commuting equations successively determine all
factor coordinates from one, and the remaining equation is commutation
with the cycle product. Therefore coset regularity is equivalent to
regularity of that product.

On Lambda\_0/pi Lambda\_0, the cycle product is the reduction of

\[
 \gamma^d/t^r=\beta=h^d.
\]

Since h lies in O{[}beta{]}, its reduction is a polynomial in the
reduction of beta. Conversely beta is h\^{}d. Thus, for the induced
m-by-m matrices,

\[
 k[\bar h]=k[\bar\beta].
 \tag{11}
\]

Their centralizers coincide. Hence beta-bar is regular if and only if
h-bar is regular. This proves that the regular points of
Y\_\{K,tau\}(gamma) correspond exactly to h-stable O\_F-lattices on
which h has regular reduction. An isomorphism of auxiliary ind-schemes
is not required: the equivariant point identification and the
transitivity argument below suffice for the component count.

For completeness, the required ordinary GL\_m transitivity has an
elementary proof. Write R = O\_F{[}h{]}. At a regular point, h-bar is a
cyclic matrix on Lambda/pi Lambda. Lift a cyclic vector to v in Lambda.
The vectors

\[
 v,hv,\ldots,h^{m-1}v
\]

are an O\_F-basis of Lambda, so Lambda = R v. Since h is regular
semisimple, any two such cyclic vectors differ by multiplication by an
element of F{[}h{]}\^{}times = A\^{}times. Thus this centralizer acts
transitively on the regular lattices. The regular locus is nonempty: the
same construction from any generic cyclic vector gives a lattice on
which h is the integral companion matrix, with cyclic reduction. The
eigenvalues of h are units, so the companion matrix is invertible over
O\_F.

Using (8)-\/-(11), the original centralizer acts transitively on the
k-valued points of Y\_\{K,tau\}\^{}\{reg\}(gamma). This suffices for
transitivity on its top-dimensional components: over the algebraically
closed field k, choose in each component a k-point belonging to no other
component, and carry one such point to the other. The spaces are locally
of finite type (or their perfections), so such points exist. Proposition
2.7 and its proof identify these components with those of
Y\_\{K,tau\}(gamma), so the latter also have one centralizer orbit.

This establishes the GL slope-block transitivity required in (1).

\section{5. Completion for GL\_n}\label{completion-for-gl_n}

Every Newton Levi of GL\_n is a direct product of general-linear groups.
On each factor the corresponding endpoint has a central Newton slope
r\_j/d\_j and length-zero rotation determined by its determinant
valuation. The Hodge-\/-Newton representative gamma\textquotesingle{}
can be taken in P\_p dot(tau), and its restriction to each factor is
isoclinic with that slope. This also follows directly from its chain
condition:
(gamma\textquotesingle\_j)\textsuperscript{\{d\_j\}/t}\{r\_j\} preserves
the relevant lattices and is a unit.

The affine Dynkin diagram is the disjoint union of the factor diagrams,
and log-volumes add. Therefore maximal log-volume in the product means
maximal log-volume in every factor. Each factor is exactly the
partial-chain case in §2. The regularity conditions and the centralizer
actions factor as well.

Section 4 gives transitivity on the regular locus of every factor. Their
product gives (1), and hence n\_\{p,gamma\} = 1 for every p in R. The
identities in §1 prove the theorem for GL\_n, including all basic and
nonbasic Newton points.

\section{6. Completion for Sp\_2n}\label{completion-for-sp_2n}

For a dominant Newton point of Sp\_2n, the positive slope blocks occur
paired with their negative duals. Its split Newton Levi has the actual
direct-product form

\[
 M\simeq\prod_{j=1}^s\mathrm{GL}_{a_j}\times\mathrm{Sp}_{2c}.
\]

The symplectic factor is the zero-slope block. A GL factor acts on a
positive-slope space and by its inverse transpose on the paired
negative-slope space; there is no determinant constraint on that factor.
Thus these are genuine GL factors, not a quotient or a subgroup with an
extra norm constraint.

The GL factors are treated by §§2-\/-4. All denominators d\_j are prime
to char(k): each a\_j \textless= n, while the paper\textquotesingle s
symplectic characteristic assumption excludes 2 and all primes
\textless= n.

For the zero-slope factor, Sp\_2c is simply connected, so its
length-zero group Omega is trivial. The tau component there is
consequently the identity. Its very special parahorics are hyperspecial.
Concretely, in affine type C\_c, deleting vertex i gives a finite Weyl
group of type C\_i times C\_\{c-i\}, with longest length

\[
 i^2+(c-i)^2.
\]

The maximal value c\^{}2 occurs exactly at i=0,c. A smaller spherical
subset cannot achieve that value. These two vertices give the
hyperspecial parahorics.

Theorem 3.5 gives centralizer transitivity on the regular locus of this
hyperspecial affine Springer fiber. The theorem applies to either
hyperspecial type by transport: conjugation by the symplectic similitude
diag(t I\_c,I\_c) gives an L-group automorphism interchanging the two
types. After moving the relevant hyperspecial subgroup to the standard
one, an element in it can be conjugated into its Iwahori by choosing a
stable isotropic flag in its reduction. This puts the situation under
the stated hypotheses of Theorem 3.5. When c=0 there is no symplectic
factor to treat.

As in the GL case, maximal log-volume, regularity, and centralizers
factor over this direct product. The general-linear transitivity from §4
and the symplectic transitivity from Theorem 3.5 therefore give (1) for
every top-contributing path. It follows that every n\_\{p,gamma\} equals
one. Section 1 proves the desired equality for Sp\_2n.

\section{7. Scope and dependency
audit}\label{scope-and-dependency-audit}

The new step is the root construction in the finite integral algebra B,
followed by the regular-locus identification (11). This handles
arbitrary slope multiplicities, rather than only the superbasic case. It
does \textbf{not} extend the local field merely to clear denominators
and then equate component counts; it identifies the original
lattice-chain problem with an F-linear lattice problem inside the
original representation, preserving its full centralizer.

The geometric inputs used from He are Theorems 1.3, 1.4, 2.8, and 3.2,
Proposition 2.5, Proposition 2.7 and its regular-locus argument, the
stabilizer argument in §3.2, and the flag/Grassmannian component
bijections in §3.6. Theorem 3.5 is needed only for the zero-slope
symplectic factor. None of these inputs is applied outside the
hypotheses stated in the uploaded paper.

The lattice calculations, integral-root existence, centralizer equality,
regularity comparison, and affine A/C maximal-log-volume classifications
are supplied in this note. Finite-rank checks of the diagram
classifications and a non-superbasic slope-block example are supporting
computations, not substitutes for the general proofs and not a Lean
verification.

The reproducible check script is \url{checks/check_local_reduction.py}.
Running it with Python in this workspace passed 77 GL rank/rotation
cases through N=12, symplectic diagram ranks 1-\/-8, the repeated GL\_4
slope 1/2 cyclic-lattice calculation, and the associated Sp\_8
block-pairing identity.

All of the local arguments use only a complete discrete valuation ring O
with algebraically closed residue field and d invertible in O. They
therefore apply to W(k) as well as k{[}{[}t{]}{]}. The auxiliary field F
may be ramified in mixed characteristic; no affine Springer transitivity
theorem over F is invoked, since the cyclic-lattice argument proves the
required transitivity directly. For the zero-slope symplectic factor the
field remains the original L, and Theorem 3.5 explicitly includes mixed
characteristic. Thus the proof covers both settings in §1.1 of the
paper.

The statement proved here is restricted to GL\_n and Sp\_2n under the
paper\textquotesingle s characteristic assumptions. Other reductive
groups, bad characteristic, and a formal proof assistant development are
outside this argument. No assertion about priority or the status of
later literature is made.
\end{document}
